\documentclass[11pt]{article}
\usepackage[a4paper,margin=24mm]{geometry}
\usepackage{amsmath,amssymb,amsthm,mathtools,bm}
\usepackage{booktabs,longtable,array}
\usepackage{enumitem}
\usepackage{xcolor}
\usepackage{microtype}
\usepackage[hidelinks]{hyperref}

\setlength{\emergencystretch}{3em}
\newtheorem{theorem}{Theorem}[section]
\newtheorem{proposition}[theorem]{Proposition}
\newtheorem{lemma}[theorem]{Lemma}
\newtheorem{corollary}[theorem]{Corollary}
\theoremstyle{definition}
\newtheorem{definition}[theorem]{Definition}
\newtheorem{gate}[theorem]{Fail-closed gate}
\theoremstyle{remark}
\newtheorem{remark}[theorem]{Remark}

\newcommand{\dd}{\mathrm d}
\newcommand{\ii}{\mathrm i}
\newcommand{\ZZ}{\mathbb Z}
\newcommand{\RR}{\mathbb R}
\newcommand{\CC}{\mathbb C}
\newcommand{\CP}{\mathbb {CP}}
\newcommand{\cO}{\mathcal O}
\newcommand{\Tr}{\operatorname{Tr}}
\newcommand{\Hol}{\operatorname{Hol}}
\newcommand{\Arf}{\operatorname{Arf}}
\newcommand{\ind}{\operatorname{ind}}
\newcommand{\rank}{\operatorname{rank}}
\newcommand{\status}[1]{\textcolor{blue!60!black}{\textsf{#1}}}
\newcolumntype{L}[1]{>{\raggedright\arraybackslash}p{#1}}

\title{AP-E3 Global-Definition and Deep-UV Audit:\\
Differential Characters on Non-Extendible Four-Sectors\\
and the Exact Boundary of the \(SU(3)\) Embedding}
\author{Route F independent research note}
\date{16 July 2026}

\begin{document}
\maketitle

\begin{abstract}
This note closes the ordinary global-definition problem for the level-two
mixed Wess--Zumino--Witten functional and sharply delimits what remains UV
data.  For the actual two-flavour target
\(X=SU(2)\times S^2\simeq S^3\times S^2\), the integral curvature
\(n\omega_3\wedge\omega_2\) has a unique degree-five differential-character
refinement.  Its holonomy defines the mixed action on every closed oriented
four-manifold and every sigma-model map, including configurations that do not
extend over a spin five-manifold.  A \v{C}ech--Deligne representative and its
descent equations make the construction local and computable.

Spin bordism contributes additional information that the curvature cannot
fix.  Stable splitting gives
\[
 \Omega^{\mathrm{Spin}}_4(S^3\times S^2)
 \simeq \ZZ\oplus\ZZ_2\oplus\ZZ_2,
\]
not merely one reduced \(\ZZ_2\).  The two reduced signs are detected by the
spin one-manifold \(g^{-1}(p_3)\) and the Arf invariant of the spin
two-manifold \(s^{-1}(p_2)\).  The full action is therefore
\[
 Z_{n,\epsilon_3,\epsilon_2}
 =\Hol_{\widehat\Omega_n}(g,s)
 (-1)^{\epsilon_3\nu_3+\epsilon_2\nu_2}.
\]
The two signs require a specified UV APS/Dai--Freed determinant and are not
silently chosen here.

The proposed all-scale breaking
\(SU(3)\to[SU(2)_c\times U(1)_g]/\ZZ_2\) is then audited exactly.  Its kernel
forces \(2j+q=0\pmod2\), so a colour singlet has even abelian charge and the
original charge-one scalar cannot occur in any \(SU(3)\) representation.
Two \(\overline{\bm3}\) scalars do furnish a conditional charge-two variant,
and explicit adjoint mass matrices can make their coloured partners heavy.
Vectorlike fundamental fermions give
\(\bm3\to\bm2_{+1}\oplus\bm1_{-2}\); all dynamical gauge anomalies cancel,
and an adjoint mass matrix can tune the quarks light while making the extra
leptons heavy.  Neither tuning is symmetry protected.  More decisively, the
charge-two scalar needs only one conjugate dressing of \(qq\), giving a
doublet rather than the exactly-two triplet.  Thus the global mixed-WZW
functional is defined, while a Route-E-equivalent all-scale UV completion is
not established.
\end{abstract}

\tableofcontents

\section{Scope, conventions, and claim discipline}

The purpose of this note is narrow.  It does not repeat the local anomaly
ledger or the worldline reduction already established in the main AP-E3
note.  Instead it addresses two residual questions:
\begin{enumerate}[label=(Q\arabic*)]
\item how to define the mixed functional when the four-dimensional spin
  configuration is not a five-dimensional spin boundary;
\item whether the suggested \(SU(3)\) embedding actually reproduces the
  charge-one scalar, decouples all partners, and preserves the level-two
  triplet mechanism.
\end{enumerate}

We use
\begin{equation}
 X=SU(2)\times S^2\simeq S^3\times S^2,
 \qquad
 \int_{S^3}\omega_3=1,
 \qquad
 \int_{S^2}\omega_2=1,
 \label{eq:normalized-forms}
\end{equation}
and the mixed curvature
\begin{equation}
 \Omega_{5,n}=n\,\omega_3\wedge\omega_2,
 \qquad n=2.
 \label{eq:mixed-curvature}
\end{equation}
The trace in
\(\omega_3=(24\pi^2)^{-1}\Tr(g^{-1}\dd g)^3\) is fixed by
Eq.~\eqref{eq:normalized-forms}, so no later sign or normalization is hidden
in a representation convention.

Statements are marked as follows.  \status{Proved} means a topological,
group-theoretic, or algebraic result follows from the displayed assumptions.
\status{Conditional} means an explicit Lagrangian or mass matrix exists but
its hierarchy or radiative stability is assumed.  \status{Open} means no
completion is claimed.

\section{The extension-free differential-character functional}

\subsection{Integral class and uniqueness}

Let \(u_3\in H^3(S^3;\ZZ)\) and
\(u_2\in H^2(S^2;\ZZ)\) be positive generators.  K\"unneth's theorem gives
\begin{equation}
 H^\bullet(X;\ZZ)
 \simeq \Lambda_{\ZZ}(u_3)\otimes
 \frac{\ZZ[u_2]}{(u_2^2)},
 \qquad
 H^4(X;\ZZ)=0,
 \qquad
 H^5(X;\ZZ)=\ZZ\{u_3u_2\}.
 \label{eq:cohomology-ring}
\end{equation}

A degree-\(k\) Cheeger--Simons differential character is a homomorphism
\begin{equation}
 \widehat h:Z_{k-1}(X)\longrightarrow\RR/\ZZ
 \label{eq:character-definition}
\end{equation}
for which there is a closed integral-period \(k\)-form
\(\operatorname{curv}(\widehat h)\) satisfying
\begin{equation}
 \widehat h(\partial C_k)
 =\int_{C_k}\operatorname{curv}(\widehat h)\pmod\ZZ.
 \label{eq:character-boundary}
\end{equation}
Differential characters are equivalent to smooth Deligne cohomology
\cite{CheegerSimons1985,Brylinski1993}.

Choose the unique normalized refinements
\begin{equation}
 \widehat u_3\in\widehat H^3(S^3;\ZZ),
 \qquad
 \widehat u_2\in\widehat H^2(S^2;\ZZ),
 \label{eq:factor-characters}
\end{equation}
with curvatures \(\omega_3\) and \(\omega_2\).  Their differential cup
product is
\begin{equation}
 \widehat\Omega_n
 :=n\,\widehat u_3\smile\widehat u_2
 \in\widehat H^5(X;\ZZ),
 \label{eq:differential-cup}
\end{equation}
and obeys
\begin{align}
 \operatorname{curv}(\widehat\Omega_n)
 &=n\omega_3\wedge\omega_2,\notag\\
 c(\widehat\Omega_n)&=n\,u_3u_2.
 \label{eq:cup-data}
\end{align}

\begin{theorem}[Unique ordinary differential refinement]
For every \(n\in\ZZ\), Eq.~\eqref{eq:differential-cup} is the unique
ordinary degree-five differential character with curvature
\(n\omega_3\wedge\omega_2\).
\end{theorem}
\begin{proof}
The curvature exact sequence is
\begin{equation}
 0\longrightarrow H^4(X;\RR/\ZZ)
 \longrightarrow\widehat H^5(X;\ZZ)
 \xrightarrow{\operatorname{curv}}\Omega^5_{\ZZ}(X)
 \longrightarrow0.
 \label{eq:curvature-sequence}
\end{equation}
Equation~\eqref{eq:cohomology-ring} and the universal coefficient theorem
give \(H^4(X;\RR/\ZZ)=0\).  Hence the fiber over the prescribed integral
curvature contains exactly one character.
\end{proof}

\subsection{Definition on an arbitrary four-sector}

Let \(M_4\) be a closed oriented four-manifold and
\(\Phi=(g,s):M_4\to X\).  Define
\begin{equation}
 Z_n^{\rm diff}[M_4,\Phi]
 :=\exp\!\left[2\pi\ii\,
 \widehat\Omega_n\bigl(\Phi_*[M_4]\bigr)\right].
 \label{eq:global-holonomy}
\end{equation}
This formula requires neither a five-manifold nor an extension of \(\Phi\).
It is therefore a definition, rather than an instruction that fails in a
non-bounding sector.

When a singular five-chain \(C_5\) in \(X\) satisfies
\(\partial C_5=\Phi_*[M_4]\), the character axiom gives
\begin{equation}
 Z_n^{\rm diff}[M_4,\Phi]
 =\exp\!\left[2\pi\ii n\int_{C_5}\omega_3\wedge\omega_2\right].
 \label{eq:chain-formula}
\end{equation}
Such a chain exists here because \(H_4(X;\ZZ)=0\).  It need not be the image
of a smooth spin five-manifold.  If a genuine extension
\((Y_5,\widetilde\Phi)\) exists, Eq.~\eqref{eq:chain-formula} reduces to
\begin{equation}
 Z_n^{\rm diff}
 =\exp\!\left[2\pi\ii n\int_{Y_5}
 \widetilde\Phi^*(\omega_3\wedge\omega_2)\right].
 \label{eq:five-manifold-limit}
\end{equation}
Changing the extension adds an integral period and leaves the exponential
unchanged.  Thus the five-dimensional formula is a consequence of the
character, not its universal definition.

\section{A computable \v{C}ech--Deligne representative}

\subsection{Descent data and triangulated holonomy}

On a good cover \(\{U_i\}\) of \(X\), a degree-five Deligne cocycle may be
written
\begin{equation}
 \check C=
 \bigl(C_i^{(4)},C_{ij}^{(3)},C_{ijk}^{(2)},
 C_{ijkl}^{(1)},C_{ijklm}^{(0)},N_{ijklmn}\bigr),
 \label{eq:deligne-data}
\end{equation}
where \(N\) is integer valued.  We fix the descent-sign convention
\begin{align}
 \dd C_i^{(4)}&=\Omega_{5,n},\notag\\
 \delta C^{(4)}&=\dd C^{(3)},&
 \delta C^{(3)}&=-\dd C^{(2)},\notag\\
 \delta C^{(2)}&=\dd C^{(1)},&
 \delta C^{(1)}&=-\dd C^{(0)},\notag\\
 \delta C^{(0)}&=N,&\delta N&=0.
 \label{eq:descent-equations}
\end{align}
The integer cocycle represents \(n u_3u_2\).  Another representative differs
by a Deligne coboundary and defines the same character.

Choose a triangulation subordinate to the cover and a barycentric labelling.
The logarithm of the holonomy is the alternating flag sum
\begin{align}
 \frac{1}{2\pi\ii}\log\Hol_{\check C}(M_4)
 ={}&\sum_{\sigma^4}\int_{\sigma^4}C^{(4)}_{i_0}
 -\sum_{\sigma^3\subset\sigma^4}\int_{\sigma^3}C^{(3)}_{i_0i_1}
 +\sum_{\sigma^2\subset\sigma^3\subset\sigma^4}
   \int_{\sigma^2}C^{(2)}_{i_0i_1i_2}\notag\\
 &-\sum_{\sigma^1\subset\cdots\subset\sigma^4}
   \int_{\sigma^1}C^{(1)}_{i_0i_1i_2i_3}
 +\sum_{\sigma^0\subset\cdots\subset\sigma^4}
   C^{(0)}_{i_0i_1i_2i_3i_4}\pmod\ZZ.
 \label{eq:cech-holonomy}
\end{align}
The orientations and incidence numbers are those of the nested simplices.
Equations~\eqref{eq:descent-equations} make internal faces cancel, while
\(N\in\ZZ\) makes cover changes and refinements invisible after
exponentiation.  This is the promised extension-free local algorithm.

\subsection{Explicit normalization check}

Use hyperspherical coordinates on \(S^3\):
\begin{equation}
 \omega_3=\frac{1}{2\pi^2}\sin^2\chi\sin\theta\,
 \dd\chi\wedge\dd\theta\wedge\dd\varphi.
 \label{eq:s3-volume}
\end{equation}
Two illustrative gerbe potentials are
\begin{align}
 B_N&=f_N(\chi)\sin\theta\,\dd\theta\wedge\dd\varphi,
 &f_N(\chi)&=\frac{2\chi-\sin2\chi}{8\pi^2},\notag\\
 B_S&=f_S(\chi)\sin\theta\,\dd\theta\wedge\dd\varphi,
 &f_S(\chi)&=f_N(\chi)-\frac1{4\pi}.
 \label{eq:gerbe-potentials}
\end{align}
They satisfy
\begin{equation}
 \dd B_N=\dd B_S=\omega_3,
 \qquad
 \int_{S^2}(B_N-B_S)=1.
 \label{eq:gerbe-overlap}
\end{equation}
The two-chart presentation may be refined to a good cover; its overlap
already exposes the integral gerbe transition.

For the \(S^2\) factor, Route-E conventions give
\begin{align}
 A_N&=\frac{1-\cos\Theta}{2}\dd\Phi,&
 A_S&=-\frac{1+\cos\Theta}{2}\dd\Phi,\notag\\
 A_N-A_S&=\dd\Phi,&
 \omega_2&=\frac{\dd A}{2\pi}
 =\frac{\sin\Theta}{4\pi}\dd\Theta\wedge\dd\Phi.
 \label{eq:hopf-patches}
\end{align}
Locally one may start with
\begin{equation}
 C_{N,\alpha}^{(4)}=n B_N\wedge\omega_2,
 \qquad
 C_{S,\alpha}^{(4)}=n B_S\wedge\omega_2.
 \label{eq:local-four-potential}
\end{equation}
The non-exact integral overlap in Eq.~\eqref{eq:gerbe-overlap} is precisely
why a single global four-form does not exist.  The remaining terms in
Eq.~\eqref{eq:deligne-data} are generated by the Deligne cup product of the
gerbe and Hopf-line cocycles.  Omitting them would make the action patch
dependent.

\section{Spin bordism, the two torsion signs, and APS selection}

\subsection{The special \texorpdfstring{\(N_f=2\)}{Nf=2} stable splitting}

The target in this repository uses \(SU(2)\simeq S^3\), so its bordism must
not be replaced by a generic \(SU(N\ge3)\) result.  The stable splitting is
\begin{equation}
 \Sigma^\infty(S^3\times S^2)_+
 \simeq S^0\vee S^2\vee S^3\vee S^5.
 \label{eq:stable-splitting}
\end{equation}
Since spin bordism is a generalized homology theory and
\begin{equation}
 \Omega^{\rm Spin}_0=\ZZ,\quad
 \Omega^{\rm Spin}_1=\ZZ_2,\quad
 \Omega^{\rm Spin}_2=\ZZ_2,\quad
 \Omega^{\rm Spin}_3=0,\quad
 \Omega^{\rm Spin}_4=\ZZ,
 \label{eq:low-spin-bordism}
\end{equation}
we obtain the following result.

\begin{theorem}[Full and reduced four-bordism]
\begin{align}
 \Omega^{\rm Spin}_4(S^3\times S^2)
 &\simeq\ZZ\oplus\ZZ_2\oplus\ZZ_2,\notag\\
 \widetilde\Omega^{\rm Spin}_4(S^3\times S^2)
 &\simeq\ZZ_2\oplus\ZZ_2.
 \label{eq:bordism-result}
\end{align}
The free \(\ZZ\) is the field-independent gravitational sector.  The first
reduced \(\ZZ_2\) comes from the \(S^3\) factor and the second from \(S^2\).
\end{theorem}

The \(S^3\) summand is the special two-flavour discrete WZW sector discussed
in Ref.~\cite{LeeOhmoriTachikawa2021Global}.  The extra \(S^2\) summand is the
one retained in the generic mixed-target calculation of
Ref.~\cite{DavighiLohitsiri2024Global}.  Both are present for the actual
product \(S^3\times S^2\).

\subsection{Transverse inverse-image invariants}

Let \(p_3\in S^3\) and \(p_2\in S^2\) be regular values.  Transversality gives
\begin{align}
 L_g&:=g^{-1}(p_3),&\dim L_g&=1,\notag\\
 \Sigma_s&:=s^{-1}(p_2),&\dim\Sigma_s&=2.
 \label{eq:inverse-images}
\end{align}
The oriented tangent spaces at the regular values frame the normal bundles.
Together with the spin structure on \(M_4\), these framings induce spin
structures on \(L_g\) and \(\Sigma_s\).  Define
\begin{align}
 \nu_3(M_4,g)&=[L_g]\in\Omega_1^{\rm Spin}\simeq\ZZ_2,
 \label{eq:nu3}\\
 \nu_2(M_4,s)&=\Arf(\Sigma_s)\in
 \Omega_2^{\rm Spin}\simeq\ZZ_2.
 \label{eq:nu2}
\end{align}
Changing the regular value gives a spin bordism between inverse images, so
the invariants are well defined.  Representatives are
\begin{align}
 &(S^1_{\rm periodic}\times S^3,
 g=\operatorname{pr}_{S^3},\ s={\rm const})
 &&\text{for }\nu_3=1,\notag\\
 &(T^2_{\rm odd}\times S^2,
 g={\rm const},\ s=\operatorname{pr}_{S^2})
 &&\text{for }\nu_2=1.
 \label{eq:bordism-generators}
\end{align}

\begin{definition}[Full spin-refined mixed functional]
After factoring out the pure gravitational phase, the most general
spin-refined action with curvature Eq.~\eqref{eq:mixed-curvature} is
\begin{equation}
 Z_{n,\epsilon_3,\epsilon_2}[M_4,g,s]
 =Z_n^{\rm diff}[M_4,g,s]
 (-1)^{\epsilon_3\nu_3(M_4,g)+
       \epsilon_2\nu_2(M_4,s)},
 \qquad \epsilon_3,\epsilon_2\in\ZZ_2.
 \label{eq:full-spin-action}
\end{equation}
\end{definition}

The ordinary differential character fixes the continuous/integral mixed
part uniquely.  It does not fix either torsion character.  On a spin-bounding
configuration both \(\nu_3\) and \(\nu_2\) vanish, so all four choices reduce
to Eq.~\eqref{eq:five-manifold-limit}; this explains why extension tests alone
cannot distinguish them.

\subsection{What APS/Dai--Freed must compute}

The coefficients \((\epsilon_3,\epsilon_2)\) are selected by the phase of a
specified UV determinant, not by choosing a convenient square root.  For a
real or pseudoreal one-parameter family of regulated boundary operators
\(\mathcal D_t\), the mod-two spectral flow satisfies
\begin{equation}
 \operatorname{SF}_2(\mathcal D_t)
 =\ind_{2,{\rm APS}}(\partial_t+\mathcal D_t).
 \label{eq:aps-spectral-flow}
\end{equation}
For a complex family, the corresponding object is the exponentiated reduced
eta invariant
\begin{equation}
 \tau_Y(\mathcal D)=
 \exp\!\left[2\pi\ii\,\bar\eta(\mathcal D_Y)\right],
 \qquad
 \bar\eta=\frac{\eta+\dim\ker\mathcal D}{2},
 \label{eq:dai-freed}
\end{equation}
viewed as a section of the determinant line.  APS gluing makes the result a
bordism character \cite{AtiyahPatodiSinger1975,DaiFreed1994}.

To determine Eq.~\eqref{eq:full-spin-action}, one must evaluate the UV
determinant on both backgrounds in Eq.~\eqref{eq:bordism-generators}.  The
standard two-flavour QCD result gives
\(\epsilon_3=n_c\bmod2=0\) for two colours, provided the usual quark
determinant and skyrmion identification are retained.  No calculation in the
present repository fixes \(\epsilon_2\).  A convention
\(\epsilon_2=0\) is allowed but is not a first-principles derivation.  Modern
bordism treatments illustrate precisely why this determinant step is
independent of local anomaly matching \cite{Freed2008,Saito2024Global}.

\begin{gate}[Global-definition result]
The ordinary non-extendible mixed functional is now defined by
Eq.~\eqref{eq:global-holonomy}.  The full spin theory remains a two-bit family
until a UV APS/Dai--Freed calculation selects \((\epsilon_3,\epsilon_2)\).
Calling the entire spin refinement \emph{closed} before that calculation
would be false.
\end{gate}

\section{The exact \texorpdfstring{\(SU(3)\)}{SU(3)} quotient and branching}

\subsection{Group homomorphism and parity theorem}

Take an adjoint field \(\Sigma\sim\bm8\) with
\begin{equation}
 \langle\Sigma\rangle=V_\Sigma T,
 \qquad T=\operatorname{diag}(1,1,-2).
 \label{eq:adjoint-vev}
\end{equation}
The unbroken subgroup is the image of
\begin{equation}
 \iota:SU(2)_c\times U(1)_g\longrightarrow SU(3),
 \qquad
 \iota(U,e^{\ii\alpha})
 =\begin{pmatrix}e^{\ii\alpha}U&0\\0&e^{-2\ii\alpha}\end{pmatrix}.
 \label{eq:embedding-map}
\end{equation}
Its kernel is
\begin{equation}
 \ker\iota=\{(\bm1,1),(-\bm1,-1)\}\simeq\ZZ_2,
 \label{eq:embedding-kernel}
\end{equation}
so the gauge group below \(V_\Sigma\) is
\begin{equation}
 U(2)\simeq[SU(2)_c\times U(1)_g]/\ZZ_2.
 \label{eq:u2-group}
\end{equation}

\begin{theorem}[Quotient-parity rule]
An \(SU(2)_c\) spin-\(j\) multiplet of integer-normalized abelian charge
\(q\) is a representation of Eq.~\eqref{eq:u2-group} if and only if
\begin{equation}
 2j+q=0\pmod2.
 \label{eq:quotient-parity}
\end{equation}
Consequently every colour doublet has odd charge and every colour singlet has
even charge.
\end{theorem}
\begin{proof}
The kernel generator acts as
\((-1)^{2j}e^{\ii\pi q}\).  It must act trivially, which is exactly
Eq.~\eqref{eq:quotient-parity}.
\end{proof}

This is a representation-independent obstruction.  In particular, a
\(\bm1_{+1}\) scalar cannot descend from any \(SU(3)\) representation.  It is
not repaired by placing that scalar next to coloured partners; the forbidden
singlet itself never appears.

\subsection{Low-representation decomposition}

With the normalization in Eq.~\eqref{eq:adjoint-vev}, direct weight
decomposition gives \cite{Slansky1981}
\begin{longtable}{@{}L{0.16\textwidth}L{0.72\textwidth}@{}}
\toprule
\(SU(3)\) irrep & \(SU(2)_c\times U(1)_g\) content\\
\midrule
\endfirsthead
\toprule
\(SU(3)\) irrep & \(SU(2)_c\times U(1)_g\) content\\
\midrule
\endhead
\(\bm3\) & \(\bm2_{+1}\oplus\bm1_{-2}\)\\
\(\overline{\bm3}\) & \(\bm2_{-1}\oplus\bm1_{+2}\)\\
\(\bm6\) & \(\bm3_{+2}\oplus\bm2_{-1}\oplus\bm1_{-4}\)\\
\(\overline{\bm6}\) & \(\bm3_{-2}\oplus\bm2_{+1}\oplus\bm1_{+4}\)\\
\(\bm8\) & \(\bm3_0\oplus\bm2_{+3}\oplus\bm2_{-3}\oplus\bm1_0\)\\
\(\bm{10}\) & \(\bm4_{+3}\oplus\bm3_0\oplus\bm2_{-3}\oplus\bm1_{-6}\)\\
\bottomrule
\end{longtable}
Every term obeys Eq.~\eqref{eq:quotient-parity}.  The companion verifier
derives the symmetric-power rows from weights rather than storing only their
dimensions.

\section{Fermions, anomaly matching, and lepton decoupling}

\subsection{Vectorlike deep-UV multiplets}

Take \(N_f=2\) Dirac fields \(\Psi_i\in\bm3\) of \(SU(3)\).  Below
\(V_\Sigma\),
\begin{equation}
 \Psi_i\longrightarrow q_i\oplus\ell_i,\qquad
 q_i\sim\bm2_{+1},\qquad
 \ell_i\sim\bm1_{-2}.
 \label{eq:quark-lepton-split}
\end{equation}
In all-left notation the fields are
\begin{center}
\begin{tabular}{@{}ccc@{}}
\toprule
field & \(SU(2)_c\) & \(U(1)_g\)\\
\midrule
\(q_L\) & \(\bm2\) & \(+1\)\\
\(q_R^c\) & \(\bm2\) & \(-1\)\\
\(\ell_L\) & \(\bm1\) & \(-2\)\\
\(\ell_R^c\) & \(\bm1\) & \(+2\)\\
\bottomrule
\end{tabular}
\end{center}
The extra lepton is not optional: it is the third component of every
fundamental \(SU(3)\) fermion.

The deep theory is vectorlike, hence
\begin{equation}
 \mathcal A_{SU(3)^3}=N_f[A(\bm3)+A(\overline{\bm3})]=0.
 \label{eq:su3-anomaly}
\end{equation}
After branching, the complete local ledger is
\begin{align}
 \mathcal A_{SU(2)_c^2U(1)_g}
 &=N_fT(\bm2)(1-1)=0,\notag\\
 \mathcal A_{U(1)_g^3}
 &=N_f\{2(1^3-1^3)+(-2)^3+2^3\}=0,\notag\\
 \mathcal A_{{\rm grav}^2U(1)_g}
 &=N_f\{2(1-1)-2+2\}=0.
 \label{eq:branched-anomalies}
\end{align}
There are \(2N_f=4\) left-Weyl colour doublets, so the colour Witten anomaly
also vanishes.

For independent quark and lepton chiral-flavour backgrounds, the
operator-valued mixed coefficients are
\begin{equation}
 \widehat\kappa_q=2(+1)=+2,\qquad
 \widehat\kappa_\ell=1(-2)=-2.
 \label{eq:kappa-cancellation}
\end{equation}
They cancel on the diagonal deep-UV flavour group, as required because the
simple \(SU(3)\) phase has no exact magnetic \(U(1)\) one-form symmetry
\cite{DavighiLohitsiri2024Global}.  For \(N_f=2\), perturbative flavour cubic
anomalies vanish, but the global Witten count is informative: the quark
flavour factor sees two copies (even), the lepton factor one copy (odd), and
the diagonal deep factor three copies (odd).  Decoupling the lepton while
keeping an exact unchanged deep chiral symmetry is therefore not innocuous.

\subsection{Adjoint mass matrix and its exact conditions}

The most direct renormalizable splitting is
\begin{equation}
 \mathcal L_{\rm mass}
 =-\overline\Psi_{L\,i}
 \bigl(m_{ij}+y_{ij}\Sigma\bigr)\Psi_R^j+{\rm h.c.},
 \label{eq:fermion-mass-lagrangian}
\end{equation}
where \(m\) and \(y\) are complex \(N_f\times N_f\) matrices.  The two
blocks are
\begin{equation}
 M_q=m+V_\Sigma y,\qquad
 M_\ell=m-2V_\Sigma y.
 \label{eq:mass-blocks}
\end{equation}

\begin{proposition}[Tree-level decoupling condition]
All two quark flavours are massless and all two leptons are massive at tree
level if
\begin{equation}
 m=-V_\Sigma y,\qquad \det y\ne0.
 \label{eq:fermion-tuning}
\end{equation}
Then
\begin{equation}
 M_q=0,\qquad
 M_\ell=-3V_\Sigma y,\qquad
 \sigma_{\min}(M_\ell)=3V_\Sigma\sigma_{\min}(y)>0.
 \label{eq:fermion-spectrum}
\end{equation}
\end{proposition}

This proves a mass-matrix identity, not natural decoupling.  Both terms in
Eq.~\eqref{eq:fermion-mass-lagrangian} connect left and right flavour groups.
If the exact deep symmetry is
\(SU(N_f)_L\times SU(N_f)_R\) and \(\Sigma\) is a flavour singlet, they are
forbidden.  If they are present, only vector flavour is exact and the quark
chiral group is accidental below the thresholds.  In the absence of a
protecting symmetry, the cancellation \(M_q=0\) must be retuned after heavy
gauge, adjoint, and lepton loops; parametrically,
\begin{equation}
 \delta M_q\sim\frac{g_3^2}{16\pi^2}\,yV_\Sigma
 \label{eq:threshold-estimate}
\end{equation}
up to representation- and scheme-dependent coefficients.  The emergence of
quark flavour, rather than its exact deep-UV realization, is required by the
generalized-symmetry no-go theorem.

\begin{gate}[Fermion decoupling]
Equations~\eqref{eq:fermion-tuning}--\eqref{eq:fermion-spectrum} are
\status{proved tree-level algebra}.  A symmetry-protected or fully retuned
quantum decoupling with the required emergent chiral theory is \status{open}.
Adding spectators changes both the flavour-Witten and determinant-phase
ledgers and must be audited as a new model.
\end{gate}

\section{Scalar embedding: the no-go and the charge-two variant}

\subsection{The charge-one model cannot be embedded}

The intermediate AP-E3 model used a colour-singlet scalar doublet
\(\phi_I\sim\bm1_{+1}\), \(I=1,2\).  The quotient-parity theorem proves that
this representation does not exist in the \(SU(3)\) embedding.  There is a
second, directly physical form of the obstruction.

Let the quark doublet have positive charge \(X_q\) and a colour-singlet scalar
have positive charge \(X_\phi\).  Equation~\eqref{eq:quotient-parity} implies
\begin{equation}
 X_q\in2\ZZ+1,\qquad X_\phi\in2\ZZ.
 \label{eq:charge-parities}
\end{equation}
A two-quark baryon dressed by \(m\) conjugate scalars is neutral only if
\begin{equation}
 2X_q-mX_\phi=0.
 \label{eq:dressing-diophantine}
\end{equation}
For exactly two dressings, \(m=2\) requires
\(X_q=X_\phi\), contradicting Eq.~\eqref{eq:charge-parities}.

\begin{theorem}[\(SU(3)\) exactly-two no-go]
No representation of
\([SU(2)_c\times U(1)_g]/\ZZ_2\) descending from \(SU(3)\) can
simultaneously contain an odd-charge colour-doublet quark, an equal-charge
colour-singlet scalar, and the neutral local operator
\(qq(\phi^\dagger)^2\).  The AP-E3 exactly-two triplet mechanism therefore
does not lift unchanged to this \(SU(3)\) embedding.
\end{theorem}

\subsection{A concrete charge-two scalar branch}

There is nevertheless a nearby, anomaly-consistent scalar spectrum.  Take two
complex fields
\begin{equation}
 \Phi_I\in\overline{\bm3},\qquad I=1,2,
 \label{eq:antitriplet-scalars}
\end{equation}
forming a global \(SU(2)_\phi\) doublet.  They decompose as
\begin{equation}
 \Phi_I\longrightarrow\chi_I\oplus\phi_I,\qquad
 \chi_I\sim\bm2_{-1},\qquad
 \phi_I\sim\bm1_{+2}.
 \label{eq:scalar-decomposition}
\end{equation}
The coloured fields \(\chi_I\) are the unavoidable partners.

The quadratic potential may be written
\begin{equation}
 V_2=
 \Phi_I^\dagger
 \bigl[(m^2)_{IJ}+\kappa_{IJ}\Sigma_{\overline3}
 +\rho_{IJ}\Sigma_{\overline3}^2\bigr]\Phi_J,\qquad
 \Sigma_{\overline3}=-\Sigma_3^T.
 \label{eq:scalar-potential}
\end{equation}
After Eq.~\eqref{eq:adjoint-vev},
\begin{align}
 M_\chi^2&=m^2-V_\Sigma\kappa+V_\Sigma^2\rho,\notag\\
 M_\phi^2&=m^2+2V_\Sigma\kappa+4V_\Sigma^2\rho.
 \label{eq:scalar-mass-blocks}
\end{align}
Exact \(SU(2)_\phi\) requires each coefficient matrix to be proportional to
\(\delta_{IJ}\).  A sufficient tree-level window is
\begin{equation}
 \lambda_{\rm eff}>0,\quad
 M_\phi^2=-\mu_\phi^2<0,\quad
 M_\chi^2+\lambda_{\chi\phi}v_\phi^2>0,\quad
 M_\chi,\ M_\ell\gg v_\phi,\Lambda_c.
 \label{eq:scalar-window}
\end{equation}
For example,
\(V_\Sigma=1,m^2=1,\kappa=-1,\rho=0\) gives
\(M_\chi^2=2\) and \(M_\phi^2=-1\) before quartic shifts.  This establishes a
nonempty tree-level parameter region.

The charge-two condensate has a \(\ZZ_2\) stabilizer in the covering
\(SU(2)\times U(1)\), but this stabilizer is exactly the identified colour
centre.  Indeed \([U,-1]=[-U,1]\) in \(U(2)\), so
\begin{equation}
 U(2)\xrightarrow{\langle\phi\rangle}SU(2)_c
 \label{eq:second-breaking}
\end{equation}
without an additional independent \(\ZZ_2\) gauge factor.  The coarse scalar
vacuum is still
\begin{equation}
 \{\phi^\dagger\phi=v_\phi^2\}/U(1)_g\simeq\CP^1.
 \label{eq:charge-two-cp1}
\end{equation}

For nonminimal scalar charge, our symbol
\(\operatorname{Vol}_{S^2}\) denotes the curvature form induced from the
covering (U(1)) orbit and is normalized by
\begin{equation}
 \int_{S^2}\operatorname{Vol}_{S^2}=X_\phi,
 \qquad
 \omega_2:=\frac{\operatorname{Vol}_{S^2}}{X_\phi},
 \qquad \int_{S^2}\omega_2=1.
 \label{eq:nonminimal-s2-normalization}
\end{equation}
Consequently, in the covering-group normalization,
\begin{equation}
 S_{\rm mix}=2\pi\hbar(n_cX_q)
 \int\omega_3\wedge\frac{\operatorname{Vol}_{S^2}}{X_\phi}
 \label{eq:charge-two-wzw}
\end{equation}
would have coefficient \(n_cX_q=2\) for
\(n_c=2,X_q=1,X_\phi=2\)
\cite{DavighiLohitsiri2024Global}.  Because the actual faithful subgroup is
\(U(2)=[SU(2)\times U(1)]/\ZZ_2\), applying this covering-
\(U(1)\) normalization to the quotient theory also requires a global match
of principal bundles and allowed fluxes.  We have not performed that bundle
match, so Eq.~\eqref{eq:charge-two-wzw} is a conditional normalization result,
not an all-bundle theorem for faithful \(U(2)\).  This topological statement
must not be confused with the local dressing representation.  Equation
\eqref{eq:dressing-diophantine} now gives
\begin{equation}
 qq\phi^\dagger:\qquad 2(1)-1(2)=0.
 \label{eq:one-dressing}
\end{equation}
Only one global \(SU(2)_\phi\) doublet is used.  The operator is a doublet,
not \(\operatorname{Sym}^2\bm2=\bm3\), and no exactly-two microscopic
derivation of the Route-E triplet survives.

\begin{gate}[Faithful-\(U(2)\) bundle normalization]
The missing step is to prove that the determinant line and its first Chern
class remain the stated integral generator when the \(U(1)\) flux is
correlated with the \(SU(2)\)-centre 't~Hooft flux by the quotient.  Until
that calculation is supplied,
\texttt{u2\_quotient\_global\_bundle\_normalization\_proven=false}.
\end{gate}

\begin{gate}[Status of the charge-two branch]
The representation, anomaly, and tree-level scalar mass algebra are
\status{conditional but explicit}.  The branch is not equivalent to the
original charge-one/exactly-two model.  It may UV-complete the normalized
mixed topological coefficient only if the faithful-\(U(2)\) bundle gate
closes, while it already fails the specific triplet mechanism; these are
logically distinct claims.
\end{gate}

\section{Running, monopoles, and the all-scale boundary}

With two Dirac \(\bm3\) fermions, one real adjoint scalar, and two complex
\(\overline{\bm3}\) scalars, the one-loop deep-UV coefficient is
\begin{align}
 b_0^{SU(3)}
 &=\frac{11}{3}C_2(\bm8)
 -\frac43N_fT(\bm3)
 -\frac16T(\bm8)
 -\frac13(2)T(\overline{\bm3})\notag\\
 &=11-\frac43-\frac12-\frac13
 =\frac{53}{6}>0.
 \label{eq:su3-beta}
\end{align}
Thus this field list removes the abelian Landau pole and remains
asymptotically free at one loop.  This does not close the generalized symmetry
at finite \(V_\Sigma\).  The breaking \(SU(3)\to U(2)\) produces finite-mass
nonabelian monopoles.  They violate the abelian Bianchi identity and hence the
continuous magnetic one-form symmetry is exact only in the decoupling limit
where those monopoles are removed.  At finite scale it is emergent, as is the
separate quark chiral symmetry after the lepton mass splitting.  This is
consistent with the no-go theorem but prevents an all-scale exact two-group
claim.

The remaining nonperturbative requirements are independent:
\begin{enumerate}
\item select a mesonic rather than diquark/Pauli--G\"ursey vacuum;
\item prove a stable \(B=1\) soliton after the charge-two deformation;
\item compute the monopole-induced violation of winding conservation;
\item evaluate the two UV APS/Dai--Freed torsion characters;
\item produce either a protected lepton-decoupling symmetry or a complete
  threshold-retuning prescription;
\item recover an exactly-two triplet by a different all-scale group or accept
  that the charge-two branch addresses only the mixed topological term.
\end{enumerate}

\section{Deterministic algebraic and numerical regression}

The companion program
\begin{center}
\texttt{route\_f/code/verify\_ap\_e3\_nonextendible\_wzw\_su3\_uv.py}
\end{center}
checks, without fitting parameters:
\begin{enumerate}
\item the full and reduced stable-splitting bordism groups;
\item Gaussian-quadrature normalization of \(\omega_3\), \(\omega_2\), and
  their product;
\item the local gerbe derivative, overlap period, Hopf transition, and
  extension phases;
\item branching of
  \(\bm3,\overline{\bm3},\bm6,\overline{\bm6},\bm8,\bm{10}\)
  from explicit weights;
\item the quotient parity and exactly-two Diophantine no-go;
\item all dynamical gauge anomalies, the colour Witten parity, and
  \(\widehat\kappa_q+\widehat\kappa_\ell=0\);
\item a non-diagonal two-flavour benchmark for
  \(M_q=0,M_\ell=-3V_\Sigma y\), the scalar split, and
  \(b_0^{SU(3)}=53/6\);
\item fail-closed booleans for torsion selection, radiative protection,
  finite-scale one-form exactness, Route equivalence, and physics promotion.
\end{enumerate}

A green mechanical run is compatible with
\texttt{full\_uv\_closed=false}.  This is deliberate: numerical verification
can confirm the displayed algebra but cannot supply missing strong dynamics
or an unspecified fermion determinant.

\section{Claim ledger and next decisive calculations}

\begin{longtable}{@{}L{0.42\textwidth}L{0.16\textwidth}L{0.32\textwidth}@{}}
\toprule
Claim & Status & Boundary\\
\midrule
\endfirsthead
\toprule
Claim & Status & Boundary\\
\midrule
\endhead
Degree-five ordinary differential character on every four-cycle &
\status{proved} & unique because \(H^4(X;\RR/\ZZ)=0\)\\
\(\Omega_4^{\rm Spin}(S^3\times S^2)=\ZZ\oplus\ZZ_2\oplus\ZZ_2\) &
\status{proved} & stable splitting; free gravitational factor separated\\
Two transverse-preimage torsion invariants & \status{proved} &
\(\nu_3\in\Omega_1^{\rm Spin}\),
\(\nu_2=\Arf\in\Omega_2^{\rm Spin}\)\\
UV values of \((\epsilon_3,\epsilon_2)\) & \status{open} &
\(\epsilon_3=0\) follows only with the standard two-colour determinant;
\(\epsilon_2\) uncomputed\\
No \(SU(3)\)-derived colour-singlet charge-one scalar & \status{proved} &
quotient parity \(2j+q=0\pmod2\)\\
No exactly-two dressing in this \(SU(3)\) embedding & \status{proved} &
\(2X_q=2X_\phi\) contradicts odd/even parity\\
Charge-two \(\overline{\bm3}\) scalar mass split & \status{conditional} &
nonempty tree-level window; thresholds and vacuum require control\\
Dynamical gauge-anomaly cancellation & \status{proved} & vectorlike
\(SU(3)\) and explicit branched ledger\\
Light quarks plus heavy leptons & \status{conditional} & exact tree identity,
no chiral protection\\
Normalized mixed level \(n=2\) in charge-two branch & \status{conditional} &
covering-group result; faithful-\(U(2)\) bundle/flux matching and local
triplet remain open\\
Route-E-equivalent all-scale UV completion & \status{open} & monopoles,
strong phase, torsion, radiative decoupling, and triplet mismatch\\
\bottomrule
\end{longtable}

The ordered next calculations are:
\begin{enumerate}
\item build the actual Euclidean UV Dirac/Yukawa operator on the two generator
  backgrounds in Eq.~\eqref{eq:bordism-generators} and compute its APS
  spectral flow or Dai--Freed eta phase;
\item perform one-loop matching of Eqs.~\eqref{eq:fermion-mass-lagrangian} and
  \eqref{eq:scalar-potential}, including all symmetry-allowed counterterms;
\item quantify finite-mass monopole events and the lifetime of the emergent
  winding charge;
\item decide explicitly whether the charge-two branch is retained only as a
  mixed-WZW UV candidate or replaced by a different group capable of the
  exactly-two triplet.
\end{enumerate}

\bibliographystyle{unsrt}
\bibliography{route_f/tex/ap_e3_nonextendible_wzw_su3_uv_audit}

\end{document}
